\documentclass[11pt]{article}
\usepackage[margin=1in]{geometry}
\usepackage{amsmath,amssymb}
\usepackage[T1]{fontenc}
\usepackage{lmodern}
\usepackage{hyperref}

\title{Literature Status for arXiv:2205.13287}
\author{arxiv\_automatic\_math run fa\_banach\_001}
\date{2026-06-20}

\begin{document}
\maketitle

\section*{Status}

\textbf{Result type:} literature\_already\_answered.

\textbf{Verdict:} a weak-star counterpart question left open in
arXiv:2205.13287 has a negative answer in later literature.

\section*{Original Question}

The source paper is:
\begin{quote}
Rainis Haller, Eve Oja, and Mart Poldvere,
\emph{Diameter two properties for spaces of Lipschitz functions},
arXiv:2205.13287.
\end{quote}

After separating several non-weak-star diameter two properties for spaces of
Lipschitz functions, the paper records that it remained open whether these
properties coincide with their weak-star versions. In particular, it asks
whether, for spaces of Lipschitz functions, \(w^\ast\)-SSD2P implies
slice-D2P.

\section*{Supporting Answer}

The supporting paper is:
\begin{quote}
Rainis Haller, Jaan Kristjan Kaasik, and Andre Ostrak,
\emph{Separating diameter two properties from their weak-star counterparts in
spaces of Lipschitz functions}, arXiv:2404.11430.
\end{quote}

The supporting paper explicitly identifies the unresolved weak-star
counterpart issue from the Haller--Oja--Poldvere paper. In Section 3 it gives
a metric space
\[
  M=\{a_1,a_2,b_k,c_k : k\in\mathbb N\}
\]
with distances \(d(a_1,b_{2k-1})=d(a_2,b_{2k})=d(c_k,b_l)=1\) for
\(k\leq l\), and distance \(2\) between different points in all remaining
cases.

For this metric space, the authors verify that \(M\) has SLTP. By the known
characterization, this yields \(w^\ast\)-SSD2P for \(\mathrm{Lip}_0(M)\).
They then define a functional
\[
  G=\frac12(m_{a_1,a_2}+F),
  \qquad
  F(f)=\lim_{\mathcal U} f(m_{b_{2n-1},b_{2n}}),
\]
where \(\mathcal U\) is a free ultrafilter, and use the slice
\(S(G,\alpha)\) to show that the diameter of this slice is strictly less
than \(2\) for suitable \(\alpha\). Thus \(\mathrm{Lip}_0(M)\) does not have
the slice-D2P.

Consequently \(w^\ast\)-SSD2P does not imply slice-D2P for spaces of
Lipschitz functions. This negatively answers the selected open question from
arXiv:2205.13287.

\section*{Scope And Limitations}

This is an already-known literature answer, not a new proof produced in this
run. It closes the weak-star-counterpart question selected from
arXiv:2205.13287, and it separates \(w^\ast\)-SSD2P from slice-D2P within the
class of spaces \(\mathrm{Lip}_0(M)\).

It does not settle the separate question in arXiv:2205.13287 asking whether
there exists a metric space \(M\) such that \(\mathrm{Lip}_0(M)\) has
slice-D2P but not D2P. The later paper arXiv:2404.11430 still lists the
distinction between \((w^\ast)\)-slice-D2P and \((w^\ast)\)-D2P in Lipschitz
function spaces as open.

\section*{Run Evidence}

The lightweight run indexes were searched for arXiv ids 2205.13287 and
2404.11430 and for the phrases weak-star counterparts, \(w^\ast\)-SSD2P, and
slice-D2P. A prior packet using arXiv:2404.11430 exists for the distinct
2018 dual-space question about \(w^\ast\)-SSD2P without SSD2P, but no packet
for this exact arXiv:2205.13287 source question was found.

\begin{thebibliography}{9}
\bibitem{HOP2022}
R. Haller, E. Oja, and M. Poldvere,
\emph{Diameter two properties for spaces of Lipschitz functions},
arXiv:2205.13287.

\bibitem{HKO2024}
R. Haller, J. K. Kaasik, and A. Ostrak,
\emph{Separating diameter two properties from their weak-star counterparts in
spaces of Lipschitz functions}, arXiv:2404.11430.
\end{thebibliography}

\end{document}
